\documentclass[10pt,a4paper]{amsart}
\usepackage{amsmath,amssymb,amsthm,graphicx}
\usepackage{verbatim}
\newenvironment{reference}{\comment}{\endcomment}
%\newenvironment{reference}{}{}
\newenvironment{slogan}{\comment}{\endcomment}
\newenvironment{history}{\comment}{\endcomment}

% For commutative diagrams we use Xy-pic
\usepackage[all]{xy}

% We use 2cell for 2-commutative diagrams.
\xyoption{2cell}
% Theorem environments.
%
\theoremstyle{plain}
\newtheorem{theorem}[subsection]{Theorem}
\newtheorem{proposition}[subsection]{Proposition}
\newtheorem{lemma}[subsection]{Lemma}

\theoremstyle{definition}
\newtheorem{definition}[subsection]{Definition}
\newtheorem{example}[subsection]{Example}
\newtheorem{exercise}[subsection]{Exercise}
\newtheorem{situation}[subsection]{Situation}

\theoremstyle{remark}
\newtheorem{remark}[subsection]{Remark}
\newtheorem{remarks}[subsection]{Remarks}

\numberwithin{equation}{subsection}

% Macros
%
\def\lim{\mathop{\mathrm{lim}}\nolimits}
\def\colim{\mathop{\mathrm{colim}}\nolimits}
\def\Spec{\mathop{\mathrm{Spec}}}
\def\Hom{\mathop{\mathrm{Hom}}\nolimits}
\def\Ext{\mathop{\mathrm{Ext}}\nolimits}
\def\SheafHom{\mathop{\mathcal{H}\!\mathit{om}}\nolimits}
\def\SheafExt{\mathop{\mathcal{E}\!\mathit{xt}}\nolimits}
\def\Sch{\mathit{Sch}}
\def\Mor{\mathop{\mathrm{Mor}}\nolimits}
\def\Ob{\mathop{\mathrm{Ob}}\nolimits}
\def\Sh{\mathop{\mathit{Sh}}\nolimits}
\def\NL{\mathop{N\!L}\nolimits}
\def\CH{\mathop{\mathrm{CH}}\nolimits}
\def\proetale{{pro\text{-}\acute{e}tale}}
\def\etale{{\acute{e}tale}}
\def\QCoh{\mathit{QCoh}}
\def\Ker{\mathop{\mathrm{Ker}}}
\def\Im{\mathop{\mathrm{Im}}}
\def\Coker{\mathop{\mathrm{Coker}}}
\def\Coim{\mathop{\mathrm{Coim}}}

% Boxtimes
%
\DeclareMathSymbol{\boxtimes}{\mathbin}{AMSa}{"02}

%
% Macros for moduli stacks/spaces
%
\def\QCohstack{\mathcal{QC}\!\mathit{oh}}
\def\Cohstack{\mathcal{C}\!\mathit{oh}}
\def\Spacesstack{\mathcal{S}\!\mathit{paces}}
\def\Quotfunctor{\mathrm{Quot}}
\def\Hilbfunctor{\mathrm{Hilb}}
\def\Curvesstack{\mathcal{C}\!\mathit{urves}}
\def\Polarizedstack{\mathcal{P}\!\mathit{olarized}}
\def\Complexesstack{\mathcal{C}\!\mathit{omplexes}}
% \Pic is the operator that assigns to X its picard group, usage \Pic(X)
% \Picardstack_{X/B} denotes the Picard stack of X over B
% \Picardfunctor_{X/B} denotes the Picard functor of X over B
\def\Pic{\mathop{\mathrm{Pic}}\nolimits}
\def\Picardstack{\mathcal{P}\!\mathit{ic}}
\def\Picardfunctor{\mathrm{Pic}}
\def\Deformationcategory{\mathcal{D}\!\mathit{ef}}
\usepackage{cleveref}
\crefname{originallemma}{original lemma}{original lemmas}
\Crefname{originallemma}{Original lemma}{Original lemmas}
\makeatletter
\newlabel{algebra-lemma-NAK}{{00DV}{source}{}{sourcepointer-1}{}}
\newlabel{algebra-lemma-NAK@cref}{{[originallemma][0][]00DV}{[1][0][]source}{}{}{}}
\makeatother
\crefname{originallemma}{original lemma}{original lemmas}
\Crefname{originallemma}{Original lemma}{Original lemmas}
\makeatletter
\newlabel{lemma-support-covering}{{07TY}{source}{}{sourcepointer-2}{}}
\newlabel{lemma-support-covering@cref}{{[originallemma][0][]07TY}{[1][0][]source}{}{}{}}
\makeatother
\crefname{originallemma}{original lemma}{original lemmas}
\Crefname{originallemma}{Original lemma}{Original lemmas}
\makeatletter
\newlabel{morphisms-lemma-support-finite-type}{{056J}{source}{}{sourcepointer-3}{}}
\newlabel{morphisms-lemma-support-finite-type@cref}{{[originallemma][0][]056J}{[1][0][]source}{}{}{}}
\makeatother
\crefname{originallemma}{original lemma}{original lemmas}
\Crefname{originallemma}{Original lemma}{Original lemmas}
\makeatletter
\newlabel{spaces-properties-lemma-stalk-pullback-quasi-coherent}{{05VQ}{source}{}{sourcepointer-4}{}}
\newlabel{spaces-properties-lemma-stalk-pullback-quasi-coherent@cref}{{[originallemma][0][]05VQ}{[1][0][]source}{}{}{}}
\makeatother
\crefname{originalsection}{original section}{original sections}
\Crefname{originalsection}{Original section}{Original sections}
\makeatletter
\newlabel{spaces-properties-section-properties-modules}{{05VR}{source}{}{sourcepointer-5}{}}
\newlabel{spaces-properties-section-properties-modules@cref}{{[originalsection][0][]05VR}{[1][0][]source}{}{}{}}
\makeatother
\title{The support of a finite type quasi-coherent module on an algebraic space is closed and computes by stalks}
\author{The Stacks Project Authors}
\date{stacks-project:spaces-morphisms.tex:local-snapshot-20261012; GNU FDL 1.2 or later}
\begin{document}
\maketitle
\setcounter{subsection}{1}
\noindent\textit{Editorial note.} All delivered mathematical text is the author's own, copied from the pinned Stacks Project chapter file spaces-morphisms.tex; no statement or proof step was rewritten, reordered or removed.\par\medskip
\noindent\textit{Editorial note.} The proof handles the three assertions in three paragraphs of one proof environment; the geometric point argument in the third paragraph is the same device as in the second.\par\medskip
\noindent\textit{Editorial note.} The five references of the counted spans are recorded as non-load-bearing navigation pointers; the argument that consumes them stays inside the counted proof span.\par\medskip
\noindent\textit{Editorial note.} Preamble: the three native excerpt spans of sources/stacks-spaces-morphisms/preamble.tex (lines 8-13, 14-18 and 154-217). The chapter file itself carries no preamble text.\par\medskip
\noindent\textit{Allowed standard prerequisites (editorial).}\begin{itemize}
\item The preimage of the support under an etale morphism is the support of the pullback. Lemma lemma-support-covering, used at proof lines 3075-3077. Reduces closedness of the support to the scheme case.
\item The support of a finite type quasi-coherent module on a scheme is closed. Morphisms, Lemma morphisms-lemma-support-finite-type, used at proof lines 3076-3077. Is the scheme input for assertion (1).
\item Nakayama's lemma. Algebra, Lemma algebra-lemma-NAK, used at proof lines 3081-3083. Gives the second equivalence in assertion (2).
\item Stalk of a pullback of a quasi-coherent sheaf. Properties of Spaces, Lemma spaces-properties-lemma-stalk-pullback-quasi-coherent, used at proof lines 3090-3097. Computes the stalk of f\textasciicircum{}*F at a geometric point as a tensor product.
\end{itemize}\medskip
\ifdefined\hypertarget\hypertarget{sourcepointer-1}{}\fi
\noindent\textit{Source reference.} Original lemma 00DV: focused-seeds/AGP/sources/stacks-snapshot-20261011/tags.txt, source lines 512--512. Navigation pointer into the Stacks Project; the standard fact it supplies is applied inside the counted proof span and the argument resting on it stays there.\par\medskip
\ifdefined\hypertarget\hypertarget{sourcepointer-2}{}\fi
\noindent\textit{Source reference.} Original lemma 07TY: focused-seeds/AGP/sources/stacks-snapshot-20261011/tags.txt, source lines 9953--9953. Navigation pointer into the Stacks Project; the standard fact it supplies is applied inside the counted proof span and the argument resting on it stays there.\par\medskip
\ifdefined\hypertarget\hypertarget{sourcepointer-3}{}\fi
\noindent\textit{Source reference.} Original lemma 056J: focused-seeds/AGP/sources/stacks-snapshot-20261011/tags.txt, source lines 6712--6712. Navigation pointer into the Stacks Project; the standard fact it supplies is applied inside the counted proof span and the argument resting on it stays there.\par\medskip
\ifdefined\hypertarget\hypertarget{sourcepointer-4}{}\fi
\noindent\textit{Source reference.} Original lemma 05VQ: focused-seeds/AGP/sources/stacks-snapshot-20261011/tags.txt, source lines 7559--7559. Navigation pointer into the Stacks Project; the standard fact it supplies is applied inside the counted proof span and the argument resting on it stays there.\par\medskip
\ifdefined\hypertarget\hypertarget{sourcepointer-5}{}\fi
\noindent\textit{Source reference.} Original section 05VR: focused-seeds/AGP/sources/stacks-snapshot-20261011/tags.txt, source lines 7560--7560. Navigation pointer into the Stacks Project; the standard fact it supplies is applied inside the counted proof span and the argument resting on it stays there.\par\medskip
\begin{lemma}
\label{lemma-support-finite-type}
Let $S$ be a scheme. Let $X$ be an algebraic space over $S$.
Let $\mathcal{F}$ be a finite type quasi-coherent $\mathcal{O}_X$-module.
Then
\begin{enumerate}
\item The support of $\mathcal{F}$ is closed.
\item For a geometric point $\overline{x}$ lying over $x \in |X|$ we have
$$
x \in \text{Supp}(\mathcal{F})
\Leftrightarrow
\mathcal{F}_{\overline{x}} \not = 0
\Leftrightarrow
\mathcal{F}_{\overline{x}} \otimes_{\mathcal{O}_{X, \overline{x}}}
\kappa(\overline{x}) \not = 0.
$$
\item For any morphism of algebraic spaces $f : Y \to X$ the pullback
$f^*\mathcal{F}$ is of finite type as well and we have
$\text{Supp}(f^*\mathcal{F}) = f^{-1}(\text{Supp}(\mathcal{F}))$.
\end{enumerate}
\end{lemma}

\begin{proof}
Choose a scheme $U$ and a surjective \'etale morphism $\varphi : U \to X$.
By Lemma \ref{lemma-support-covering} the inverse image of the support of
$\mathcal{F}$ is the support of $\varphi^*\mathcal{F}$ which is closed by
Morphisms, Lemma \ref{morphisms-lemma-support-finite-type}.
Thus (1) follows from the definition of the topology on $|X|$.

\medskip\noindent
The first equivalence in (2) is the definition of support.
The second equivalence follows from Nakayama's lemma, see
Algebra, Lemma \ref{algebra-lemma-NAK}.

\medskip\noindent
Let $f : Y \to X$ be as in (3). Note that $f^*\mathcal{F}$ is of finite type
by Properties of Spaces, Section
\ref{spaces-properties-section-properties-modules}.
For the final assertion, let $\overline{y}$ be a geometric point of $Y$
mapping to the geometric point $\overline{x}$ on $X$. Recall that
$$
(f^*\mathcal{F})_{\overline{y}} =
\mathcal{F}_{\overline{x}} \otimes_{\mathcal{O}_{X, \overline{x}}}
\mathcal{O}_{Y, \overline{y}},
$$
see Properties of Spaces, Lemma
\ref{spaces-properties-lemma-stalk-pullback-quasi-coherent}.
Hence $(f^*\mathcal{F})_{\overline{y}} \otimes \kappa(\overline{y})$
is nonzero if and only if
$\mathcal{F}_{\overline{x}} \otimes \kappa(\overline{x})$ is nonzero.
By (2) this implies $x \in \text{Supp}(\mathcal{F})$ if and only
if $y \in \text{Supp}(f^*\mathcal{F})$, which is the content of
assertion (3).
\end{proof}

\end{document}
